\section{答题时间规划与答题技巧}

\subsection{试卷结构速览}

\begin{center}
\begin{tabular}{c|c|c|c}
\textbf{题型} & \textbf{题量} & \textbf{分值} & \textbf{建议用时} \\ \toprule
听力 & 20 题 & 30 分 & 20 分钟 \\
阅读理解（A--D） & 15 题 & 37.5 分 & 35 分钟 \\
七选五 & 5 题 & 12.5 分 & 10 分钟 \\
完形填空 & 15 题 & 15 分 & 15 分钟 \\
语法填空 & 10 题 & 15 分 & 10 分钟 \\
短文改错 & 10 题 & 10 分 & 10 分钟 \\
书面表达 & 1 题 & 25 分 & 20 分钟 \\ \hline
\textbf{总计} & \textbf{76 题} & \textbf{150 分} & \textbf{120 分钟}
\end{tabular}
\end{center}

\subsection{答题顺序策略}

\textbf{推荐顺序}（听力固定在前，其余灵活调整）：

\begin{enumerate}
  \item \textbf{听力}（20 分钟）：注意力高度集中，利用读题时间预判
  \item \textbf{阅读理解 A-B 篇}（15 分钟）：相对简单，确保全对
  \item \textbf{七选五}（10 分钟）：重逻辑，不宜恋战
  \item \textbf{完形填空}（15 分钟）：先通读再选，注意上下文线索
  \item \textbf{阅读理解 C-D 篇}（20 分钟）：难度较大，留够时间
  \item \textbf{语法填空 + 短文改错}（20 分钟）：基础题型，细心拿分
  \item \textbf{书面表达}（20 分钟）：先列提纲再写，留检查时间
\end{enumerate}

\textbf{底线}：基础题（听力+阅读AB+七选五+语填+改错，约 80 分）应确保 90\% 以上得分率。

\subsection{听力技巧}

\begin{enumerate}
  \item \textbf{预读题目}：充分利用试音和每段间隔读题，划出关键词
  \item \textbf{关注信号词}：but, however, actually, in fact 后常是答案
  \item \textbf{数字题}：注意区分 -teen / -ty，记录数字做简单计算
  \item \textbf{推断题}：对话隐含信息，不选原文原词
  \item \textbf{独白题}：注意首句（主旨）、列举（细节）、转折（重点）
  \item \textbf{速记技巧}：符号+缩写记关键信息
\end{enumerate}

\subsection{阅读理解（四选一）技巧}

\begin{enumerate}
  \item \textbf{先题后文}：先读题目（不看选项），划关键词，再读文章定位
  \item \textbf{主旨大意题}（ what's the main idea / best title）：
    \begin{itemize}
      \item 关注首段、尾段、每段首句
      \item 排除过于宽泛或过于狭小的选项
      \item 标题类注意概括性和吸引力
    \end{itemize}
  \item \textbf{细节理解题}（what / when / where / why / how）：
    \begin{itemize}
      \item 定位到原文后，同义替换是常见命题方式
      \item 绝对词（all, never, always）常为错误选项
    \end{itemize}
  \item \textbf{推理判断题}（infer / imply / suggest / conclude）：
    \begin{itemize}
      \item 选的不是原文直接说的，而是推测出的
      \item 排除与原文矛盾或无关的选项
    \end{itemize}
  \item \textbf{词义猜测题}（"..." probably means）：
    \begin{itemize}
      \item 利用上下文、构词法、同义反义关系推测
      \item 代入验证法
    \end{itemize}
  \item \textbf{作者态度题}（attitude / tone）：
    \begin{itemize}
      \item 关注形容词、副词、情态动词
      \item 常见态度：positive, negative, neutral, critical, supportive, doubtful
    \end{itemize}
\end{enumerate}

\subsection{七选五技巧}

\begin{enumerate}
  \item \textbf{先定首尾}：首句常是主题句，尾句常是结论
  \item \textbf{看逻辑词}：however（转折）、therefore（因果）、for example（举例）、also（并列）
  \item \textbf{代词指代}：this, that, it, they 等指代前文内容
  \item \textbf{段落结构}：总--分结构，注意空格在段首/段中/段尾
  \item \textbf{排除法}：先确定有明显线索的选项，逐步缩小范围
  \item \textbf{带入验证}：选完所有空后通读一遍，确保行文连贯
\end{enumerate}

\subsection{完形填空技巧}

\begin{enumerate}
  \item \textbf{先通读}：不看选项通读全文，理解大意
  \item \textbf{上下文线索}（最核心）：
    \begin{itemize}
      \item 前文提示：空格答案在前面出现过
      \item 后文提示：读完后半段才能确定
      \item 逻辑关系：and, but, so, however, otherwise
    \end{itemize}
  \item \textbf{词义辨析}：
    \begin{itemize}
      \item 动词辨析：注意动作先后、施动者
      \item 名词辨析：注意可数不可数、抽象具体
      \item 形容词辨析：注意褒贬色彩
    \end{itemize}
  \item \textbf{固定搭配}：动词短语、介词搭配需平时积累
  \item \textbf{复现原则}：同一词或同义词在上下文中重复出现
  \item \textbf{验证}：选完后代入，看语义是否通顺
\end{enumerate}

\subsection{语法填空技巧}

\begin{enumerate}
  \item \textbf{有提示词}（7 题）：
    \begin{itemize}
      \item 动词：先判断谓语/非谓语
      \begin{itemize}
        \item 谓语 $\rightarrow$ 时态、语态、主谓一致
        \item 非谓语 $\rightarrow$ to do / doing / done
      \end{itemize}
      \item 名词：单复数、所有格
      \item 形容词/副词：比较级、最高级、词性转换
      \item 代词：主宾格、物主代词、反身代词
    \end{itemize}
  \item \textbf{无提示词}（3 题）：
    \begin{itemize}
      \item 介词：固定搭配、时间/地点介词
      \item 冠词：a / an / the，注意泛指和特指
      \item 连词：and, but, or, so, when, while, because, if
      \item 从句引导词：which, that, who, whose, where, when, why
    \end{itemize}
  \item \textbf{注意大小写}：句首单词首字母大写
  \item \textbf{检查}：通读一遍，看语法是否合理
\end{enumerate}

\subsection{短文改错技巧}

\begin{enumerate}
  \item \textbf{考点分布规律}（10 个错，各 1 分）：
    \begin{itemize}
      \item 动词时态/语态 1--2 处
      \item 名词单复数 1 处
      \item 冠词 1 处
      \item 形容词/副词混淆 1 处
      \item 介词 1 处
      \item 连词 1 处
      \item 代词 1 处
      \item 非谓语动词 1 处
      \item 固定搭配/词形 1 处
    \end{itemize}
  \item \textbf{修改方式}：
    \begin{itemize}
      \item 删除（多一词）$\rightarrow$ 划线
      \item 添加（漏一词）$\rightarrow$ 加插入符号
      \item 修改（错一词）$\rightarrow$ 划线改正
    \end{itemize}
  \item \textbf{通读法}：逐句读，第一遍找明显错误，第二遍找细微错误
  \item \textbf{常见陷阱}：
    \begin{itemize}
      \item 动词加-s 问题（主谓一致）
      \item 及物/不及物动词
      \item 可数/不可数名词
      \item 比较级/最高级
    \end{itemize}
\end{enumerate}

\subsection{书面表达技巧}

\begin{enumerate}
  \item \textbf{时间分配}：3 分钟审题列提纲，15 分钟写作，2 分钟检查
  \item \textbf{审题要点}：
    \begin{itemize}
      \item 文体：书信/通知/续写/议论文/记叙文
      \item 人称：第一/第二/第三人称
      \item 时态：一般现在/一般过去/一般将来
      \item 要点：逐条覆盖，不要遗漏
    \end{itemize}
  \item \textbf{结构框架}：
    \begin{itemize}
      \item 开头：开门见山，说明写作目的
      \item 正文：2--3 段，每段一个要点
      \item 结尾：总结/期待回应
    \end{itemize}
  \item \textbf{高级表达}：
    \begin{itemize}
      \item 句式多样：复合句、倒装句、强调句、非谓语
      \item 词汇替换：important $\to$ significant, many $\to$ numerous, because $\to$ due to
      \item 逻辑连接：firstly, moreover, furthermore, in addition, consequently
    \end{itemize}
  \item \textbf{扣分点}：
    \begin{itemize}
      \item 要点遗漏（最致命）
      \item 时态混乱
      \item 拼写错误
      \item 字数不足/超出
    \end{itemize}
\end{enumerate}

\subsection{时间管理技巧}

\begin{itemize}
  \item \textbf{听力不分心}：听力一旦开始，注意力必须集中，错过不纠结
  \item \textbf{阅读卡壳}：阅读题思考超过 2 分钟无进展，先跳下一题
  \item \textbf{完形不恋战}：第一遍做不出来的空先跳过，第二遍缩小范围
  \item \textbf{写作留检查}：务必留 2 分钟检查拼写、时态、单复数
  \item \textbf{涂卡时间}：建议做完一题涂一题，或每大题结束后集中涂
  \item \textbf{最后 5 分钟}：检查答题卡填涂、未做的题选最可能答案
\end{itemize}
